\chapterpage{Matrices}{Six problems. Two-dimensional arrays are mostly an
index-discipline test: get \texttt{rows} and \texttt{cols} the right way round
and most of the difficulty disappears.}

% ------------------------------------------------------------------
\problem{Rotate a square matrix 90 degrees clockwise}{Medium}

Rotate in place. \texttt{[[1,2],[3,4]]} becomes \texttt{[[3,1],[4,2]]}.

\idea{Transpose (reflect across the main diagonal), then reverse each row. Two
simple passes beat one confusing index formula.}

\solheading
\begin{lstlisting}[style=java]
public static void rotate90(int[][] m) {
    int n = m.length;
    for (int i = 0; i < n; i++)
        for (int j = i + 1; j < n; j++) {      // j starts at i+1
            int t = m[i][j]; m[i][j] = m[j][i]; m[j][i] = t;
        }
    for (int[] row : m)
        for (int i = 0, j = n - 1; i < j; i++, j--) {
            int t = row[i]; row[i] = row[j]; row[j] = t;
        }
}
\end{lstlisting}

\cost{$O(n^2)$ time, $O(1)$ space.}
\trap{Starting the inner transpose loop at \texttt{j = 0} swaps every pair
twice and leaves the matrix unchanged. It must start at \texttt{i + 1}.}

% ------------------------------------------------------------------
\problem{Spiral traversal}{Medium}

Return the elements in spiral order: left to right along the top, down the
right side, right to left along the bottom, up the left side, then inwards.

\idea{Four moving boundaries. After walking each edge, shrink that boundary, and
check before the two ``return'' edges that the boundaries have not crossed.}

\solheading
\begin{lstlisting}[style=java]
public static List<Integer> spiral(int[][] m) {
    List<Integer> out = new ArrayList<>();
    if (m == null || m.length == 0) return out;
    int top = 0, bottom = m.length - 1;
    int left = 0, right = m[0].length - 1;

    while (top <= bottom && left <= right) {
        for (int c = left; c <= right; c++) out.add(m[top][c]);
        top++;
        for (int r = top; r <= bottom; r++) out.add(m[r][right]);
        right--;
        if (top <= bottom) {
            for (int c = right; c >= left; c--) out.add(m[bottom][c]);
            bottom--;
        }
        if (left <= right) {
            for (int r = bottom; r >= top; r--) out.add(m[r][left]);
            left++;
        }
    }
    return out;
}
\end{lstlisting}

\cost{$O(\text{rows}\times\text{cols})$ time.}
\trap{Omitting the two guards duplicates a row or a column on non-square
matrices --- a single remaining row would be walked both left-to-right and
right-to-left.}

% ------------------------------------------------------------------
\problem{Count islands}{Medium}

The grid holds \texttt{'1'} for land and \texttt{'0'} for water. Count the
connected groups of land, using four-way adjacency.

\idea{Scan every cell. The first time you touch a piece of land, increment the
count and flood the whole island so it is never counted again.}

\solheading
\begin{lstlisting}[style=java]
public static int countIslands(char[][] grid) {
    if (grid == null || grid.length == 0) return 0;
    int rows = grid.length, cols = grid[0].length, islands = 0;
    int[][] dirs = {{-1,0},{1,0},{0,-1},{0,1}};

    for (int r = 0; r < rows; r++)
        for (int c = 0; c < cols; c++) {
            if (grid[r][c] != '1') continue;
            islands++;
            Deque<int[]> queue = new ArrayDeque<>();
            queue.add(new int[] {r, c});
            grid[r][c] = '0';                   // mark on push
            while (!queue.isEmpty()) {
                int[] cell = queue.poll();
                for (int[] d : dirs) {
                    int nr = cell[0] + d[0], nc = cell[1] + d[1];
                    if (nr < 0 || nr >= rows || nc < 0 || nc >= cols) continue;
                    if (grid[nr][nc] != '1') continue;
                    grid[nr][nc] = '0';
                    queue.add(new int[] {nr, nc});
                }
            }
        }
    return islands;
}
\end{lstlisting}

\cost{$O(\text{rows}\times\text{cols})$ time and space.}
\trap{Marking a cell visited when you \emph{pop} it rather than when you
\emph{push} it lets the same cell enter the queue many times, and a large
open grid then blows up. Mark on push, always.}

% ------------------------------------------------------------------
\problem{Set matrix zeroes}{Medium}

If any cell is zero, set its whole row and column to zero. The catch is doing
it without a second matrix.

\idea{Two passes. First record \emph{which} rows and columns must be cleared;
only then clear them. Clearing as you go turns the whole matrix to zero.}

\solheading
\begin{lstlisting}[style=java]
public static void setZeroes(int[][] m) {
    if (m == null || m.length == 0) return;
    int rows = m.length, cols = m[0].length;
    boolean[] clearRow = new boolean[rows];
    boolean[] clearCol = new boolean[cols];

    for (int r = 0; r < rows; r++)
        for (int c = 0; c < cols; c++)
            if (m[r][c] == 0) { clearRow[r] = true; clearCol[c] = true; }

    for (int r = 0; r < rows; r++)
        for (int c = 0; c < cols; c++)
            if (clearRow[r] || clearCol[c]) m[r][c] = 0;
}
\end{lstlisting}

\cost{$O(\text{rows}\times\text{cols})$ time, $O(\text{rows}+\text{cols})$
space.}
\trap{The naive one-pass version writes zeros that the rest of the pass then
reads as original zeros, cascading until everything is zero. The $O(1)$-space
variant stores the flags in row 0 and column 0 and needs a separate flag for
each of those --- more clever than it is worth under time pressure.}

% ------------------------------------------------------------------
\problem{Diagonal sums}{Easy}

Sum the main diagonal and the anti-diagonal of a square matrix, counting the
centre of an odd-sized matrix once.

\idea{One loop. At step $i$ the two diagonal cells are $m[i][i]$ and
$m[i][n-1-i]$; they coincide exactly when $i = n-1-i$.}

\solheading
\begin{lstlisting}[style=java]
public static long diagonalSum(int[][] m) {
    int n = m.length;
    long total = 0;
    for (int i = 0; i < n; i++) {
        total += m[i][i];
        if (i != n - 1 - i) total += m[i][n - 1 - i];
    }
    return total;
}
\end{lstlisting}

\cost{$O(n)$ time, $O(1)$ space.}
\trap{Double-counting the centre cell of an odd matrix. And the total can
exceed \texttt{int} --- use \texttt{long}.}

% ------------------------------------------------------------------
\problem{Search a row-and-column sorted matrix}{Medium}

Every row is sorted left to right and every column top to bottom. Decide
whether a target is present.

\idea{Start at the top-right corner. That cell is the largest in its row and the
smallest in its column, so a comparison always eliminates an entire row or an
entire column.}

\solheading
\begin{lstlisting}[style=java]
public static boolean contains(int[][] m, int target) {
    if (m == null || m.length == 0) return false;
    int r = 0, c = m[0].length - 1;
    while (r < m.length && c >= 0) {
        if (m[r][c] == target) return true;
        if (m[r][c] > target) c--;     // whole column is too big
        else r++;                      // whole row is too small
    }
    return false;
}
\end{lstlisting}

\cost{$O(\text{rows}+\text{cols})$ time, $O(1)$ space.}
\trap{Starting at the top-left corner gives no information: both moves lead to
larger values, so you cannot eliminate anything. The corner choice is the
entire trick.}
